\section*{\centering ВСТУП}\addcontentsline{toc}{section}{ВСТУП}

\textbf{Обґрунтування вибору теми дослідження.} Стрімкий розвиток технологій штучного інтелекту (ШІ) докорінно змінює професійну діяльність дизайнерів. Аналіз публікаційної динаміки засвідчує 6,4-кратне зростання досліджень генеративного ШІ в дизайн-освіті за останні роки, що свідчить про формування нової парадигми~--- перехід від <<дизайнера-виробника>> до <<дизайнера-куратора>>, який оцінює, редагує та спрямовує результати, згенеровані ШІ. За таких умов інформаційно-комунікативна компетентність (ІКК) стає ключовою складовою професійної підготовки майбутніх дизайнерів.

Проте аналіз 122 українських освітніх програм спеціальності 022 Дизайн, проведений із застосуванням обчислювальних методів (тематичне моделювання, мережевий аналіз, аналіз розривів), виявив системну прогалину: 0\% покриття компонентів <<Основи ШІ>>, <<Генеративний ШІ>> та <<Комп'ютерний зір>> при загальному рівні інтеграції ШІ-тематики лише 35,7\%. Водночас опитування випускників та роботодавців трьох програм КДПУ засвідчило незадоволеність рівнем ІКТ-підготовки (33\% у графічному дизайні) та прямі рекомендації роботодавців щодо інтеграції ШІ у навчальний процес.

Актуальність дослідження зумовлена суперечностями між:
\begin{itemize}
    \item стрімкою ШІ-трансформацією дизайнерської практики (6,4-кратне зростання публікацій) та відставанням програм професійної підготовки;
    \item структурною відповідністю стандарту В2 Європейській рамці кваліфікацій та 0\% покриттям компонентів ШІ у 122 українських програмах дизайну;
    \item потребою випускників і роботодавців у розвитку ІКТ-навичок та недостатнім рівнем формування ІКК у чинних навчальних програмах.
\end{itemize}

\textbf{Зв'язок роботи з науковими програмами, планами, темами.} Дисертацію виконано у межах наукового дослідження <<Формування інформаційно-комунікативної компетентності у фаховій підготовці майбутніх дизайнерів>> (державний реєстраційний номер ХХХХХХХХХХХХ). Тему дисертації затверджено вченою радою Криворізького державного педагогічного університету.

\textbf{Об'єкт дослідження}~-- процес фахової підготовки майбутніх дизайнерів у закладах вищої освіти.

\textbf{Предмет дослідження}~-- формування інформаційно-комунікативної компетентності майбутніх дизайнерів.

\textbf{Мета дослідження}~-- теоретично обґрунтувати та розробити модель формування інформаційно-комунікативної компетентності в процесі фахової підготовки майбутніх дизайнерів на основі даних, включаючи конкретні рекомендації щодо трансформації навчальних програм.

\textbf{Гіпотеза дослідження.} Формування ІКК майбутніх дизайнерів може бути системно вдосконалено, якщо: (1)~визначити теоретичні засади ІКК у контексті ШІ-трансформації дизайн-освіти; (2)~провести комплексний аналіз поточного стану ІКК-спрямованого змісту в українських освітніх програмах дизайну із застосуванням обчислювальних методів; (3)~розробити модель трансформації навчальних програм на основі даних, що визначає інтеграцію ШІ-компетентностей для кожного типу навчальних дисциплін.

\textbf{Завдання дослідження:}
\begin{enumerate}
    \item Визначити теоретичні засади формування ІКК у контексті трансформації дизайн-освіти ери ШІ.
    \item Провести комплексний аналіз поточного стану формування ІКК в українській дизайн-освіті на основі даних (загальнонаціональний рівень + кейс КДПУ).
    \item Розробити та обґрунтувати модель трансформації навчальних програм для формування ІКК з конкретними рекомендаціями на рівні окремих дисциплін.
\end{enumerate}

\textbf{Методи дослідження.} Для розв'язання поставлених завдань використано комплекс методів у межах змішаного дослідницького дизайну:
\begin{itemize}
    \item \textit{теоретичні}: аналіз, синтез, систематизація наукових джерел; оглядове дослідження (scoping review); порівняльний аналіз стандартів та освітніх програм; моделювання;
    \item \textit{емпіричні}: аналіз документів (122 програми та 134 силабуси), анкетування (32 випускники, 9 роботодавців), експертне оцінювання, контент-аналіз відгуків стейкхолдерів та акредитаційних звітів;
    \item \textit{обчислювальні}: тематичне моделювання (BERTopic), мережевий аналіз навчальних програм, аналіз ШІ-розривів, семантична подібність;
    \item \textit{статистичні}: описова статистика, критерій Краскела--Уолліса, крос-табуляції.
\end{itemize}

\textbf{Експериментальна база дослідження.} Дослідження здійснювалося на базі Криворізького державного педагогічного університету (кафедра дизайну) із залученням даних 122 акредитованих програм спеціальності 022~Дизайн (галузь знань 02~Культура і мистецтво), отриманих з Єдиної державної електронної бази з питань освіти та сайту НАЗЯВО.

\textbf{Наукова новизна одержаних результатів:}

\textit{уперше:}
\begin{itemize}
    \item операціоналізовано ІКК для дизайнерів у контексті ШІ з 4-компонентною структурою (інформаційно-аналітичний, комунікативний, технологічний, рефлексивний);
    \item проведено комплексний обчислювальний аналіз усіх 122 українських програм дизайну на предмет покриття ІКК (тематичне моделювання, мережевий аналіз, аналіз розривів у загальнонаціональному масштабі);
    \item розроблено модель трансформації навчальних програм на основі даних, що визначає інтеграцію ШІ для кожного типу навчальних дисциплін;
\end{itemize}

\textit{удосконалено:}
\begin{itemize}
    \item аналіз освітніх програм дизайну за допомогою обчислювальних методів (тематичне моделювання, мережевий аналіз, аналіз розривів у загальнонаціональному масштабі);
    \item розуміння розвитку компетентностей в контексті українсько-європейської інтеграції;
\end{itemize}

\textit{набуло подальшого розвитку:}
\begin{itemize}
    \item інтеграція інструментів ШІ у фахову підготовку з дизайну;
    \item підходи до модернізації навчальних програм на основі аналізу програм великого масштабу.
\end{itemize}

\textbf{Практичне значення одержаних результатів} полягає в розробці:
\begin{itemize}
    \item рекомендацій щодо трансформації навчальних програм на основі даних для всіх українських програм дизайну;
    \item конкретних пропозицій щодо модифікації дисциплін трьох програм дизайну КДПУ;
    \item моделі ШІ-інтегрованого бакалаврського навчального плану на 240~кредитів ЄКТС;
    \item аналітичного інструментарію (Python), придатного для аналізу інших спеціальностей.
\end{itemize}

\textbf{Особистий внесок здобувача.} У працях, опублікованих у співавторстві, здобувачу належать: систематичний огляд 156 досліджень генеративного ШІ в дизайн-освіті; аналіз стандарту В2 Дизайн у контексті Європейської рамки кваліфікацій; розробка та реалізація скрепера даних НАЗЯВО; обчислювальний аналіз 122 програм; розробка моделі трансформації навчальних програм.

\textbf{Апробація результатів дослідження.} Основні положення та результати дослідження доповідалися та обговорювалися на наукових конференціях:
\begin{itemize}
    \item IV Міжнародна науково-практична конференція <<Topical Issues of Modern Science, Society and Education>> (Харків, 1--3~листопада 2021~р.);
    \item VІІ Всеукраїнська науково-практична конференція <<Формула творчості: теорія і методика мистецької освіти>> (Полтава, 26~квітня 2023~р.);
    \item Всеукраїнська науково-практична конференція <<Актуальні проблеми формування естетичної культури учнівської та студентської молоді в закладах освіти України>> (Кривий Ріг, 27~квітня 2023~р.);
    \item І~Міжнародна науково-практична конференція Таврійського національного університету до 160-ї річниці від дня народження В.\,І.~Вернадського (Київ, 16--17~березня 2023~р.);
    \item Щорічна науково-практична конференція <<Образотворче мистецтво як інструмент міжкультурної комунікації в умовах глобалізації>> (Кривий Ріг, 24~лютого 2026~р.)
\end{itemize}

\textbf{Публікації.} Основні результати дисертації відображено у 7~наукових працях, у тому числі: 2~статті у наукових фахових виданнях, включених до міжнародних наукометричних баз, та 5~тез доповідей у збірниках матеріалів наукових конференцій.

\textbf{Структура та обсяг дисертації.} Дисертація складається з анотації, вступу, трьох розділів, висновків, списку використаних джерел (\arabic{citenum@totc}~найменувань) та додатків. Загальний обсяг роботи~-- \pageref{last_page}~сторінок, з них основного тексту~-- \pageref{end_main_text}~сторінок. Робота містить \totalfigures~рисунків та \totaltables~таблиць.
